\documentclass[10pt,a4paper]{amsart}
\usepackage[margin=19mm]{geometry}
\usepackage{amsmath,amssymb,mathtools}
\usepackage[scr=rsfs,cal=euler]{mathalfa}
\usepackage{xfrac}
\usepackage[unicode,hidelinks,bookmarks=false]{hyperref}
\newcommand{\ie}{i.e.\ }
\newcommand{\cf}{cf.\ }
\newcommand{\resp}{respectively\ }
\newcommand{\Subsection}{\subsection}
\providecommand{\nobreakdash}{\nobreak\hbox{-}}
\setlength{\emergencystretch}{2em}
\multlinegap0pt
\theoremstyle{plain}
\newtheorem{theorem}{Theorem}
\newtheorem{lemma}[theorem]{Lemma}
\newtheorem{corollary}[theorem]{Corollary}
\newtheorem{proposition}[theorem]{Proposition}
\newtheorem{claim}[theorem]{Claim}
\theoremstyle{definition}
\newtheorem{definition}{Definition}
\newtheorem{example}[theorem]{Example}
\title{Affine transverse foliations in sphere bundles}
\author{Ilya Alekseev, Ivan Nasonov and Gaiane Panina}
\date{Source: arXiv:2603.01079v4 (CC BY 4.0); the authors' own native \LaTeX{} source, set here in the editorial AMS \texttt{amsart} wrapper (the source's own class is \texttt{amsart}, at 12pt).}
\begin{document}
\maketitle
\noindent\textit{Editorial scope.} The counted claim of this unit is the article's Lemma 14 (source label \texttt{lemmaNonAm}), the growth-ratio bound for balls in the Cayley graph, delivered together with the authors' complete original proof. The statement and the proof are the only retained passages, and every one of their characters is the authors' own; the proof is delivered as the single primary proof body exactly as the authors' proof environment encloses it, including the commented-out block that the authors left inside it, which prints nothing. Printed numbering: the authors' own declarations of the theorem-like environments are carried unchanged, with the lemma sharing the counter of the theorem environment as in the authors' preamble, and that counter is advanced so that the delivered PDF prints the counted claim as Lemma~14, exactly as the published article does. Omitted: the rest of the article, that is its title block, abstract, introduction, the random-quasisection section that precedes the lemma and the sections that follow it; every passage retained in the published article only as surrounding context for this lemma refers to figures of the article or to earlier lemmas whose own proofs refer to those figures, so the retained unit is the maximal source-local passage of the article that is closed under its own cross-references. Nothing of the counted statement or of the counted proof is omitted. Class substitution: the source document class is the AMS \texttt{amsart} class at 12pt with the authors' own packages; the wrapper is AMS \texttt{amsart} at the editorial settings used throughout this corpus, carrying the authors' own macro and theorem-environment declarations. Reference handling: no retained passage contains a citation; the article's bibliography is part of the authors' own TeX and is bundled unchanged inside the source file, and no reference entry is transcribed, delivered or invented. No mathematics is written, repaired or re-derived here.

\setcounter{theorem}{13}
    \begin{lemma}\label{lemmaNonAm} For $R>2$ it holds that in $Cayley(\mathcal{G})$:
    $$\dfrac{|B_R| - |B_{R-1}|}{|B_{R-1}|} \leq 2\lambda-1.$$ 
    
    \end{lemma}

    \begin{proof}
    Consider  arbitrary $g\in B_{R-1}\setminus B_{R-2}$. There is at least one edge $[g  g_0] \in Cayley(\mathcal{G})$ such that $g_0 \in B_{R-2}$ --- it is an ``ancestor'' of $g$. 
    
    %It claimed also that there are at least $n$ other edges $[g g_1],...,[g g_n]$ such that $g_i \in B_{R-1}$. 

   \iffalse Indeed, consider some connected component of $g_0F\ \cap gF$, it is an $(n-1)$-dimensional. It have at least $n$ hyperfaces, \textcolor{red}{((we need some extra conditions for choosing of the fundamental domain!!!!!))} which correspond to fundamental domains, neighbor both to $g$ and $g_0$.
    \fi
   

    Hence, at most $|\mathcal{G}| -1$ edges of $Cayley(\mathcal{G})$ connect $g$ with vertices in $B_R\setminus B_{R-1}$. 
    
    So, $$\dfrac{|B_R| - |B_{R-1}|}{|B_{R-1}|}\leq |\mathcal{G}|-1 \leq 2\lambda-1.$$
    
    \end{proof}

\end{document}
